\documentclass[10pt]{article}
\usepackage{template}
\usepackage{amsmath,amssymb,microtype,tikz,graphicx}
\usetikzlibrary{arrows.meta,positioning,shapes.geometric}
\setlength{\emergencystretch}{2em}
\SetKeys{date=2026-10-19,vol=006,title={每日有趣新闻｜政策、科学与数学},author={ChatGPT 编辑部},note={\small 2026 年 10 月 19 日·第 006 期｜重点核查 10 月 18 日新闻；涉及 10 月 15—16 日的原始公告，均单独标明日期。预印本不视为已获独立证明。}}
\renewcommand{\newsitem}[3]{\par\medskip{\large\bfseries\color{newsblue}#1\par}{\footnotesize\color{newsgray}#2\par}\smallskip\color{black}#3\par\medskip{\color{newsgray}\rule{\linewidth}{0.25pt}}\par}
\newcommand{\newslead}[3]{\par\medskip{\LARGE\bfseries\color{newsblue}#1\par}\vspace{1mm}{\small\color{newsgray}#2\par}\smallskip\color{black}#3\par\medskip{\color{newsgray}\rule{\linewidth}{0.45pt}}\par}
\begin{document}
\begin{documents}
\newslead{技术出口管制的边界：从光伏产线到人工智能模型，中国新目录究竟管什么？}{时政与产业头条｜10 月 16 日两部门公告；10 月 18 日国际媒体进一步报道}{
十月十六日，中国商务部、科技部发布第 48 号公告，就修订《中国禁止出口限制出口技术目录》公开征求意见，并同步启动部分技术出口前的临时登记。十月十八日，路透社把焦点放在十二项光伏制造技术上，强调潜在限制可能影响外国企业建立本地生产能力。但只看“禁止出口光伏技术”的标题，会混淆三件性质不同的事情：拟禁止的技术、拟限制的技术，以及已经生效的临时登记要求。此次目录修订涉及人工智能、生物医药、光伏三个领域共四十项技术，其中十四项拟禁止出口，二十六项拟限制出口；正式目录尚在征求意见，截止日期为十一月十五日。\cite{MOFCOMNotice,MOFCOMBrief,ReutersSolar}

\textbf{管制对象首先是技术，不是所有使用该技术制造的商品。}例如，拟纳入管理的光伏条目涵盖多晶硅、硅片和电池制造，具体到一定技术指标的钙钛矿及叠层电池组件生产工艺，以及多晶硅工程技术。光伏制造链可粗略分成原料制备、硅片加工、电池结构与组件封装四个阶段。控制一条产线所依赖的工艺知识和关键设备技术，与限制成品太阳能板跨境销售并不是同一个法律动作。商务部答记者问明确表示，使用相关技术制造的产品不在此次拟列管的技术范围内；部分设备只有在承载达到指标的关键技术时才涉及相应管理。\cite{MOFCOMBrief}

\begin{center}
\begin{tikzpicture}[font=\scriptsize, x=1cm,y=1cm]
\draw[rounded corners=2pt,fill=blue!3,draw=blue!20] (0,0) rectangle (7.3,2.8);
\node[font=\bfseries,anchor=west] at (.20,2.53) {拟增列技术按领域分布（共 40 项）};
\fill[newsblue!85] (.25,1.88) rectangle (4.14,2.17);
\fill[orange!75] (4.14,1.88) rectangle (5.01,2.17);
\fill[teal!70] (5.01,1.88) rectangle (7.05,2.17);
\node[anchor=west] at (.24,1.59) {人工智能 23 项};
\node[anchor=west] at (2.63,1.59) {生物医药 5 项};
\node[anchor=west] at (5.00,1.59) {光伏 12 项};
\draw[newsgray,thin] (.23,1.35) -- (7.06,1.35);
\node[anchor=west,font=\bfseries] at (.25,1.07) {拟禁止 14 项 \quad 拟限制 26 项};
\node[anchor=west] at (.25,.69) {征求意见截至 2026-11-15};
\node[anchor=west] at (.25,.37) {已生效：公告所列技术出口前临时登记};
\end{tikzpicture}

{\scriptsize\color{newsgray}图 1｜编辑部依据商务部、科技部 2026 年第 48 号公告绘制。横向色带长度与各领域项目数量成比例；禁止／限制总数另列，不能据此推断每个领域内部的比例。}
\end{center}

为什么连钙钛矿也受到关注？传统晶硅电池依靠半导体结把光能转换为电能；钙钛矿材料可以在不同波段吸收光子，将其与硅电池叠层，有机会利用单一吸收层未能充分利用的太阳光谱。高效率并不只来自材料配方，薄膜均匀性、大面积涂布、封装寿命和工艺良率都可能决定实验室器件能否进入工厂。因而，一项“制造技术”的商业价值常常不体现在专利名称，而体现在成套参数、设备配置和工程经验中。路透社同时指出，TOPCon、HJT 与 XBC 等先进晶硅电池技术属于拟限制而非一概禁止的范围。\cite{ReutersSolar,PVMag}

人工智能部分的政策边界更加容易被误读。官方说明称，拟增列大模型、智能体、算力集群及信息处理相关二十三项技术；其中包括面向敏感用途的特定能力，而不是宣布所有开源模型或所有模型 API 一律不得向海外提供。已经通过开源或论文发表进入公共领域的技术，拟不纳入管理范围。对于限制类模型，海外本地化部署可能涉及技术转移；但在境内运行模型、向海外提供 API、网页或 SaaS 调用，官方答复认为通常不属于此次所说的大模型技术出口。这个区分反映了“转移模型技术”与“提供模型服务”在制度上的不同。\cite{MOFCOMBrief,APChina}

\textbf{已生效的部分，是登记，而不是全部最终限制。}公告规定，自十月十六日起至正式修订目录实施前，拟出口所列技术的经营者应按要求进行登记，涉及技术所有权转让、授权境外使用、实际控制权变更等情形。登记制度不等于自动获得最终出口许可，也不意味着每一种服务调用都需要登记。企业仍须对照具体条目和技术指标判断适用性；报道者同样不能把征求意见稿当作已经生效的永久禁令。\cite{MOFCOMNotice}

对国际产业链而言，问题不仅是“能否买到电池”，还包括海外工厂能否获得工艺包、产线调试和持续升级支持。中国政府将此解释为国家安全、技术安全与正常合作之间的管理安排，并强调不针对特定国家；国外产业观察者则关注限制对美国、印度等地本土化光伏生产计划的影响。前者是政策制定方的说明，后者是市场影响分析，两者不能互相替代。最终条目、豁免机制、许可证实际审批情况和企业调整方案，才是后续值得检验的事实。\cite{MOFCOMBrief,ReutersSolar}
}

\newslead{“能源停火”与不断升级的袭击：俄乌战场上的三套相互冲突的叙述}{国际时政深度｜10 月 18 日路透社、美联社报道；部分数字来自交战方声明}{
十月十八日，乌克兰总统泽连斯基表示，自美国总统特朗普一周前宣布针对能源设施的停火以来，俄罗斯已七十六次袭击乌克兰能源、石油和天然气部门。他特别提到文尼察、日托米尔、赫梅利尼茨基和敖德萨地区的变电站受到打击。路透社同时指出，基辅和莫斯科均未正式确认特朗普所宣布的停火协议。这一事实差异十分关键：一位第三方领导人宣称达成了安排，并不自动意味着两个交战方已就范围、起止时间、监督机制和违规认定形成共同文本。七十六次袭击的统计是乌方公布的数字，报道并未独立逐一核验。\cite{ReutersEnergy,APDrone}

同一天，战场另一端也发生大规模无人机行动。美联社报道，乌克兰无人机袭击导致圣彼得堡的石油设施起火，俄国防部称其击落了一千零十五架乌克兰无人机；俄罗斯地方官员通报圣彼得堡袭击造成一人死亡、五人受伤。这里“一千零十五架”是俄方声称击落的数量，不能被改写为独立确认的乌方发射总数。俄军同时继续攻击乌克兰能源设施，乌方地方消息称部分地区停电。战时双方都会发布有利于自身叙事的战果和损失数字，可靠报道必须保留归属，不能用一个未经交叉核实的整数制造精确感。\cite{APDrone}

\textbf{为什么能源设施成为双方重点？}能源基础设施同时具有经济和军事意义：电力影响铁路运输、通信、工业生产和医院运行，燃料供应关系到运输、农业机械与军队后勤。攻击炼油厂、油库和变电站，可能试图削弱对方的长期作战能力，也可能产生大范围民生后果。判断具体打击是否符合国际人道法，需要考察目标的军事用途、预期平民损害、可行替代方案及比例原则，不能仅凭“能源”二字推定每次行动的合法性或违法性。

此次军事行动还与华盛顿的制裁争论交织。美联社报道，特朗普十月十八日解释暂缓新增对俄制裁的决定，理由是担心措施破坏谈判；泽连斯基则认为放松施压会增强俄方继续作战的能力。美国政府强调谈判空间，乌克兰政府强调施压与安全保证，双方在政策优先顺序上存在可观察的分歧。但从公开表态推断任何一方“真正想要战争”或“必然准备放弃盟友”，都超出了现有证据。\cite{APDrone}

\textbf{无人机事件也跨越了边境。}乌克兰外交部就一架无人机进入立陶宛领空致歉；乌方发言人称其可能受俄罗斯电子战干扰而偏航。立陶宛媒体援引危机管理部门称，一辆行驶中的汽车受损、无人受伤。这里“可能受到干扰”是解释性假设，尚不能视为已完成的技术归因。它提醒我们，远程武器的导航、识别和电子对抗，会让局部军事行动产生跨境风险。\cite{APDrone}

下一步最需要的证据，是各方是否公布停火文本、是否建立独立核验机制，以及袭击地点、受损设备和民生影响能否得到可靠交叉核查。外交声明、交战方战果与可观测损害应分成三类记录；只有这样，公众才有可能判断所谓停火是有效约束、尚未落实的倡议，还是双方根本未共同接受的政治表述。
}

\newsitem{数学拾趣｜七百余篇机器生成的数学论文，如何判断其中哪些真的有价值？}{数学与形式化证明｜10 月 12 日 arXiv 预印本；10 月 19 日专题解读}{
十月十二日，研究者 Ying Xie 向 arXiv 提交预印本《What the Machine Invented》，尝试从 OpenAI 十月公开的机器生成数学稿件中提取可复用的引理、构造与证明方法。原论文描述初期目录含七百二十二份手稿、三百七十二个问题家族；但仓库目前的 README 显示七百一十九份手稿，说明目录可能在版本维护中调整。两组数字对应不同时间的快照，不能不加说明地混写为同一统计。OpenAI 仓库还明确承认部分未经形式化的结果可能有问题。\cite{XieSurvey,OpenAIMath}

一个证明看起来漂亮，不等于它是正确的；一个正确证明，也不等于其方法新颖。假设我们想证明一类图总有某种染色。作者可以先提出一个引理，把原问题化为更小的图，再通过归纳完成证明。即使最终结论后来被发现已有经典证明，这个新的缩减引理仍可能独立用于其他问题。反过来，若引理的适用条件被写错，后续推导即使每一步形式上流畅，整体也可能不成立。

因此，评价一份机器生成数学论文至少有三道不同的门槛：\textbf{正确性}要求每个关键推理成立；\textbf{新颖性}要求与已有文献比对，确认不是已知工具的改写；\textbf{可迁移性}则要看某个引理能否被其他问题使用。Xie 的调查研究三十八个问题家族，归纳了跨主题的证明机制，并将新颖性判断明确标注为暂定。论文也承认没有为这项调查实际编译和核验 Lean 形式化代码，因此它不是对原始七百余篇论文的集体认证。\cite{XieSurvey}

Lean 等证明助手的优势，是把推理落实为能够被内核检查的形式对象；它不能自动回答研究问题是否重要、命名是否恰当或文献优先权归属。形式化本身也需要准确表达假设：若编码的命题比原论文弱，机器核验通过也不意味着原论文全部主张得到证明。未来真正有用的数学新闻，不应只报“AI 解决了多少猜想”，而应展示一个具体定理、它依赖的关键引理、独立复核状况，以及它与旧理论之间的新联系。
}

\newsitem{把碰撞事件当作训练材料：CERN 与 Mistral 的科学基础模型计划}{科学与 AI｜合作声明发表于 10 月 15 日；本期作方法解读}{
欧洲核子研究中心 CERN 与 Mistral 在十月十五日宣布签署意向声明，计划探索用于高能物理的基础模型。它不是已经发布的成品模型，也不是宣称 AI 已经独立发现新粒子。合作设想是利用大量模拟粒子碰撞训练通用表示，使模型参与探测器数据重建、粒子识别和其他分析任务。与语言模型主要处理词元不同，碰撞数据包括探测器不同层的能量沉积、径迹和时间信息，其几何结构和物理约束更为明确。\cite{CERNMistral,ReutersCERN}

\begin{center}
\begin{tikzpicture}[font=\scriptsize,>=Latex,node distance=0.22cm]
\node[draw=blue!45,rounded corners=2pt,fill=blue!5,text width=1.58cm,align=center,minimum height=1.05cm] (a) {粒子碰撞\\（模拟或实测）};
\node[draw=blue!45,rounded corners=2pt,fill=blue!5,text width=1.58cm,align=center,minimum height=1.05cm,right=of a] (b) {探测器响应\\（能量、径迹）};
\node[draw=orange!70,rounded corners=2pt,fill=orange!8,text width=1.58cm,align=center,minimum height=1.05cm,right=of b] (c) {模型重建\\（识别、估计）};
\node[draw=teal!70,rounded corners=2pt,fill=teal!8,text width=1.58cm,align=center,minimum height=1.05cm,right=of c] (d) {物理分析\\（误差、检验）};
\draw[->,thick] (a) -- (b);\draw[->,thick] (b) -- (c);\draw[->,thick] (c) -- (d);
\end{tikzpicture}

{\scriptsize\color{newsgray}图 2｜编辑部原创流程示意；箭头表示分析依赖关系，并非 CERN 已部署的系统结构。模拟训练与真实数据验证必须区分。}
\end{center}

大型强子对撞机中的两个粒子束相撞后，实验人员看不到“粒子碰撞的照片”；探测器记录的是各类传感器的响应，研究者再依据磁场、几何与粒子相互作用规律重建运动轨迹。传统流程常由多个专门算法完成不同环节。基础模型方案希望学习可在多项任务之间复用的表示，但训练数据来自模拟时，就必须警惕模拟与真实探测器之间的分布偏差。一个模型在模拟数据上准确，不意味着它在未知噪声、探测器老化或罕见事件上仍可靠。

从科研方法看，真正值得追踪的是能否提供可复现的基准、物理约束测试、系统误差估计与罕见信号的假阳性分析，而非参数量大小。CERN 公告没有给出明确交付时间或最终预算；因此把合作意向写成已实现的技术突破是不准确的。\cite{CERNMistral}
}

\newsitem{星系团碰撞留下的“高音谱号”：一张多波段天文图如何讲述两个阶段？}{天文学与科学图像｜Chandra 原始图文发布于 10 月 14 日；近期研究解读}{
NASA 钱德拉 X 射线天文台十月十四日公布了一幅编号 VLSS J0318.9+5755 的星系团图像，画面中的弯曲射电结构因形似高音谱号而得名。图像并不是人眼直接看见的彩色照片：X 射线数据以蓝色显示，射电信号以红色呈现，红外观测则为背景星系提供信息。不同颜色代表不同波段和处理流程，不能直接理解为气体的真实肉眼颜色。\cite{ChandraImage}

研究者提出了两个相继发生的过程：首先，超大质量黑洞活动将高能粒子和磁场注入周围空间，产生射电辐射；之后，两个星系团发生碰撞，使原本衰减的电子群受到压缩和再加速，从而重新增强射电信号。这样的解释将黑洞反馈与大尺度星系团动力学联系起来，但仍需用射电谱形、气体温度和空间分布等证据检验。钱德拉图文给出，目标距离约十五点五亿光年，所展示区域横跨约六百万光年；图像结合了四月的 X 射线观测和其他望远镜数据。\cite{ChandraImage}

\textbf{为什么本期不直接转载那张合成照片？}钱德拉图像使用政策对其自制材料开放使用，但这幅多波段合成图还包含来自其他机构的观测数据，使用时需要分别核对权利与署名。我们因此在此提供原始发布链接，而不把未经完成版权核查的合成图重新打包进仓库。配图应当增加可核实的信息，而不是绕过图像授权条件。
}

\newsitem{也门北部空袭与沙特边境坠机：战事升级时，事实如何在相互指责中消失？}{中东局势重点报道｜10 月 18 日美联社报道}{
十月十八日，也门北部萨达省发生一次致命空袭。美联社援引胡塞武装控制地区的卫生部门称，萨哈尔地区一处民宅遭袭，至少九人死亡，其中包括五名儿童、两名妇女；另有十九人受伤，其中十一名儿童。胡塞方面指责沙特阿拉伯发动袭击，但报道发布时沙特政府尚未回应。这组伤亡数字来自当地实际控制方，具有重要的人道主义警示意义，但具体责任归属仍待独立调查。报道不能因为一个来源来自交战方就直接否认平民伤亡，也不能在缺少可核验资料时把指控写成已经裁定的事实。\cite{APMideast}

同一天，沙特一架武装直升机在靠近也门边境的纳季兰地区坠毁，两名机组人员死亡。胡塞方面声称将其击落；沙特军方发言人则称事故由技术故障导致。两种说法并存，说明仅凭事故发生与地理位置无法判断致因。若要确认是防空武器、机械故障还是操作失误，需要残骸损伤、雷达记录、维护日志和可信的调查报告。军事冲突中对战果的宣传有强烈动机，单方认领不能代替物证。\cite{APMideast}

埃及国家防务委员会也在当天宣布愿向沙特提供支持，以维护其安全稳定，同时表示将继续寻求外交降温。埃及的政治表态不等于已经派遣部队或介入作战；支援的具体方式、规模与法律依据仍有待公布。也门战争曾经历较长时间的相对停火，但当前冲突升级再次使红海航运、边境居民和国内流离失所者承受风险。\cite{APMideast}

\textbf{人道主义影响不止于前线。}战争可能阻断粮食和药品运输，损坏供水与电力系统，迫使医疗机构在物资短缺下收治伤员。判断局势不能只看哪一方占据某个城镇，更要追踪医院能否运转、救援车队是否获准通行、儿童营养状况是否恶化。国际人道法对交战方保护平民的义务并不因其政治立场不同而消失。对于十月十八日的空袭，最需要的是独立核实受害者身份、攻击手段、目标选择与责任链条，而不是从伤亡数字直接推出整场战争的全部因果解释。

区域外交也面临现实困难：沙特担忧本土机场和能源设施受到威胁，胡塞武装则继续依靠跨境打击增加谈判筹码；双方的军事行动又可能牵动伊朗、美国及周边国家的政策。任何“局势即将全面失控”或“停火很快恢复”的确定性预言都缺乏充分证据。可以确定的是，军事升级会扩大平民与区域基础设施面临的风险。\cite{APMideast}
}

\newsitem{加沙粮食援助缩减：从受助人数到实际食物，数字揭示了什么？}{人道主义与公共政策｜世界粮食计划署负责人 10 月 18 日声明}{
美联社十月十八日报道，联合国世界粮食计划署代理执行主任 Carl Skau 表示，由于资金困难，加沙接受援助的人数由约一百六十万减少至约一百一十万，减少约五十万人。该机构还称，食品配给减半，七万五千个家庭的现金援助减少，每日供应的热餐数量从约二十五万份下降到十五万份。这些是不同的统计指标：受助人数是覆盖范围，热餐份数是每日供应流量，配给减半则涉及每名受助者获得的食物数量。不能简单把它们相加，也不能据此直接计算营养不良人数。\cite{APMideast}

在持续受冲突影响的地区，粮食援助体系的脆弱性并不只表现为缺钱。食品必须经过采购、运输、仓储、分发和领取多个环节；安全状况、道路通行、燃料供应和市场价格都可能使名义预算无法转化为实际餐食。即便物资进入地区，如果配送网络遭到破坏，偏远或行动不便的人群仍可能被遗漏。现金援助在市场仍能运转时给予家庭一定选择权，但若食品供应不足、价格剧烈上涨，现金本身无法创造粮食。

援助削减尤其值得与即将到来的冬季联系起来。受损住房、饮水设施和卫生系统可能同时增加疾病与营养风险，而儿童、孕产妇、老人和慢性病患者承受的影响并不相同。对政策讨论而言，除了要求更多捐款，还应核对捐助承诺与实际到账的差异、援助准入限制、分发效率和独立监测。现有报道足以确认世界粮食计划署公开警告了援助规模缩减，却不足以把每个家庭的具体健康结果归因于单一因素。\cite{APMideast}
}

\newsitem{法国抗议的两个现场：民生诉求与暴力破坏不能混为一谈}{国际社会重点报道｜10 月 18 日路透社、美联社报道}{
十月十七日，法国各地围绕生活成本和学校条件举行大规模示威，十月十八日进入损失清点与政治回应阶段。美联社援引法国当局估计，前一天参与游行者约五十五万人，组织者则估计为八十万人；两组数字来自不同统计方法，不能不加说明地挑选一个当作精确人数。报道指出至少二千四百人被逮捕，多座城市出现抢掠、纵火和警民冲突，利摩日市政厅遭严重火灾破坏。与此同时，许多参与者和平表达了对物价、公共服务和教育资源的忧虑。\cite{APFrance,ReutersFrance}

马克龙在十月十八日电视讲话中谴责暴力，要求调查警务和情报准备是否存在不足，并表示地方政府必要时可限制部分集会。他同时要求政府邀请学生、教师、议员、工会和企业代表商议具体对策。政府承认存在需要回应的社会诉求，并不自动说明其对抗议处理的每一项措施都已得到充分监督；反过来，示威活动中发生破坏，也不能抹去和平参与者的公共政策主张。

事件还出现有关执法手段的争议。美联社报道，法国检方正在调查两名年轻人在抗议期间遭受严重伤害的事件，包括一名十七岁少女失去一只眼睛，另有一名年轻示威者失去一只手。调查启动不等于已经确定警员刑事责任，但严重伤害足以要求独立保存现场录像、执法记录和医疗证据。对警方而言，维护公共安全与保障和平集会自由不是二选一，而是需要同时履行的义务。\cite{APFrance}

不同政治阵营也对暴力起因提出互相冲突的解释。政府将部分破坏归咎于激进团体和犯罪分子，反对派则批评政府治理和警务部署。美国政治人物还把骚乱归因于移民，但美联社指出法国当局没有建立这项关联。新闻报道不能把未经证实的政治归因转化为既定事实。接下来真正值得检验的是政府是否提出可量化的民生政策、警方与司法调查是否公开、损毁设施如何修复，以及新的集会限制是否有明确法律依据和救济途径。\cite{APFrance,ReutersFrance}
}

\newsitem{厄尔尼诺可能在年末增强：气候预测为何不能直接变成灾害预言？}{气候科学专题｜世界气象组织 10 月 8 日季节展望}{
世界气象组织十月八日公布的全球季节性气候展望指出，热带太平洋中东部海表温度异常正在增强，预计厄尔尼诺在十月至十二月期间继续发展，并可能于十二月附近达到峰值。厄尔尼诺是热带太平洋海气耦合系统的异常状态，常伴随信风、海表温度和对流位置的共同变化。它会改变大尺度环流，使不同地区的降雨与温度分布偏离常年水平，但影响具有显著区域差异。\cite{WMOClimate}

这里必须区分三个层次。首先，海表温度观测可以显示已经发生的异常；其次，季节预测使用多模式集合估计未来数月可能的演变；最后，某个城市是否发生洪水或干旱，还取决于当地天气、地形、水文设施和脆弱性。全球尺度的厄尔尼诺预测不能替代具体地区的短期天气预报，更不能推出某个灾害必然在某日发生。世界气象组织强调其预报旨在帮助农业、供水、公共卫生和防灾部门提前准备，而不是给出单点确定性结论。

对新闻编辑而言，气候科学图表尤其需要保留参考时段、温度距平基准、预测集合与不确定性范围。只画一条不断上升的红线，却不说明它是观测值还是模型预测，会让读者误以为所有未来变化都已确定。本刊将继续跟踪实际观测与后续季节预测的偏差，避免把“概率升高”写成“灾难已被预言”。\cite{WMOClimate}
}

\newsitem{从星云图像到声音：NASA 再次展示科学可视化的边界}{科学传播简讯｜NASA 每日天文图片 10 月 18 日}{
NASA 十月十八日的每日天文图片展示鹰状星云“创生之柱”的声音化作品：制作团队把观测图像中的位置和亮度等视觉信息映射到音高与声音事件，使观众以听觉体验星云结构。这不是太空中真实传播到地球的声音；星际空间极其稀薄，不能把数据映射得到的音效当作录音。科学声音化可以为视觉障碍者提供额外的信息通道，也有助于公众发现图像中的模式，但映射规则必须公开，避免把人为选择的声学效果误认为天体物理量的直接测量。\cite{NASAAPOD}
}

\par\medskip{\large\bfseries\color{newsblue}参考资料与原始出处}\par
\begingroup\footnotesize\renewcommand{\refname}{}\bibliographystyle{unsrt}\bibliography{refs}\endgroup
\end{documents}
\end{document}